\section*{Övningar}
\addcontentsline{toc}{section}{Övningar}

Lösningarna till uppgifterna som inte ber om kod står i \hyperref[ch:losningar]{appendix~A}.
Koden publiceras i kursrepot efter föreläsningen.

\exercise{Hitta framen}{Avkodning}
\label{ex:2:1}
En nod tar emot följande byte-ström:

\begin{console}
32 74 00 FF A5 F7 03 03 25 17 7F 05 10 20 30 02 C2
\end{console}

\begin{enumerate}[a)]
  \item Vid vilket index börjar framen?
  \item Vilka bytes tillhör den, från SOF till och med CHK?
  \item Vad ska parsern göra med bytesen före?
  \item Vilken typ har framen, och vilka adresser går den mellan?
\end{enumerate}

\exercise{SOF-fällan}{Resonemang}
\label{ex:2:2}
Parsern står i \code{WaitForSof2} och tar emot \code{0xA5}.
\begin{enumerate}[a)]
  \item Vad gör den naiva implementationen, och vad går förlorat?
  \item Vad ska den göra i stället?
  \item Konstruera en byte-ström där skillnaden avgör om en giltig frame hittas eller inte.
\end{enumerate}

\exercise{Tappad byte}{Avkodning}
\label{ex:2:3}
Samma frame som i övning \ref{ex:2:1}, men byten \code{0x20} tappas på vägen.
\begin{enumerate}[a)]
  \item Vilket fält läser parsern som checksumma?
  \item Kommer \code{isFrameReady()} att bli sann?
  \item Vad returnerar \code{extractFrame()}, och varför?
  \item I vilket tillstånd står parsern efteråt, och vad krävs för att den ska hitta nästa frame?
\end{enumerate}

\exercise{Längdfältet}{Resonemang}
\label{ex:2:4}
En parser kontrollerar inte längdfältet mot buffertens storlek.
\begin{enumerate}[a)]
  \item Konstruera en byte-ström som får den att skriva utanför sin buffert.
  \item Varför räcker det inte att kontrollera längden i \code{extractFrame()}?
  \item Vilken allmän regel följer av det, och var i \kapref{ch:driver} gäller samma regel igen?
\end{enumerate}
